% chapters/01/section-01.tex
\section{Ruang Lingkup dan Capaian Mata Kuliah}
\label{sec:ruang-lingkup}

Permainan digital yang tampak sederhana menyimpan satu kenyataan yang jarang
disadari pemainnya. Di balik setiap lompatan karakter, setiap musuh yang
mengejar, dan setiap efek suara ketika koin terambil, ada program yang terus
berjalan dan mengambil keputusan puluhan kali setiap detik. Mata kuliah
\NamaMK\ mengajak Anda membangun program semacam itu, bukan sekadar
memainkannya.

\subsection{Makna Memprogram Permainan}

Memprogram permainan berarti menulis program yang menerima masukan pemain,
memperbarui keadaan dunia permainan, lalu menampilkan hasilnya kembali kepada
pemain, dan mengulangi ketiga langkah tersebut terus-menerus selama permainan
berjalan \cite{nystrom-gameloop}. Pola kerja ini berbeda dari program yang
biasa Anda tulis pada mata kuliah pemrograman dasar. Program penghitung nilai
rata-rata berjalan sekali, mengeluarkan hasil, lalu berhenti. Program
permainan tidak pernah berhenti sendiri; ia berhenti hanya ketika pemain
memilih keluar atau permainan mencapai keadaan akhir.

Perbedaan itu membawa konsekuensi teknis yang menjadi inti mata kuliah ini.
Karena program berjalan terus, waktu menjadi unsur yang harus dikelola, bukan
diabaikan. Karena pemain dapat menekan tombol kapan saja, masukan harus
ditangani sebagai peristiwa, bukan sebagai urutan pertanyaan yang diajukan
program. Karena objek bergerak, posisi harus diperbarui dengan memperhatikan
selang waktu antar-bingkai agar kecepatan gerak tidak bergantung pada
komputer yang dipakai. Karena objek dapat bertabrakan, perlu aturan yang
menentukan kapan dua objek dianggap bersentuhan dan apa yang terjadi
sesudahnya.

Seluruh persoalan tersebut akan Anda selesaikan sendiri pada bagian pertama
buku ini, dengan menulis kodenya dalam bahasa Python dan pustaka Pygame
\cite{sweigart-mg}. Baru setelah mekanisme dasarnya dipahami, pekerjaan
dipindahkan ke \istilah{game engine} Unity, tempat sebagian besar persoalan
tersebut sudah ditangani oleh perkakas sehingga perhatian dapat dialihkan
kepada desain level, penyajian aset, dan pengujian hasil \cite{unity-leveldesign}.

\subsection{Kedudukan Mata Kuliah dan Pengetahuan Prasyarat}

Mata kuliah ini berbobot \BobotSKS, diselenggarakan pada semester
\SemesterMK, dan menuntut mata kuliah prasyarat \MataKuliahPrasyarat.
Prasyarat tersebut bukan formalitas. Buku ini mengandaikan Anda sudah mampu
menggunakan variabel dan tipe data, menulis percabangan dan perulangan,
menyusun fungsi beserta parameternya, serta menggunakan struktur data dasar
seperti daftar dan kamus. Kemampuan membaca pesan kesalahan dan menelusuri
penyebab kegagalan program juga diandaikan sudah terbentuk.

Materi prasyarat tersebut tidak diulang dalam buku ini. Ketika Bab IV
memperkenalkan \istilah{game loop} berbentuk perulangan tanpa akhir,
pembahasannya berpijak pada pemahaman Anda tentang perulangan, bukan
mengajarkannya kembali. Jika ada bagian prasyarat yang sudah lama tidak Anda
pakai, sebaiknya disegarkan lebih dahulu; Bab II memuat pengingat singkat
mengenai konsep yang akan dipakai berulang kali.

Capaian mata kuliah ini juga menjadi bekal bagi mata kuliah lanjutan dan
tugas akhir yang menuntut pembangunan perangkat lunak interaktif. Alasannya,
kompetensi yang dilatih di sini---mengelola siklus kerja yang berjalan terus,
menangani masukan, dan menguji hasil secara sistematis---berlaku umum pada
perangkat lunak yang berinteraksi langsung dengan penggunanya.

\subsection{Ruang Lingkup dan Batasan Pembahasan}

Agar jelas apa yang menjadi tanggung jawab buku ini, batasan berikut
ditetapkan sejak awal.

\begin{itemize}
  \item \textbf{Sasaran utama adalah permainan dua dimensi} yang dibangun
        dengan Python/Pygame dan dengan Unity. Permainan tiga dimensi dibahas
        hanya sebatas konsep yang dipakai bersama, misalnya perbedaan
        \istilah{Collider2D} dan \istilah{Collider3D} pada Bab XIII.
  \item \textbf{Mekanisme internal permainan ditulis sendiri} pada Blok I dan
        Blok II, tanpa memanfaatkan \istilah{engine}. Tujuannya agar Anda
        mengenal cara kerja sistem permainan, bukan hanya cara memanggil
        perkakas.
  \item \textbf{Kecerdasan buatan yang dibahas terbatas pada AI dasar}: NPC
        yang mengejar pemain, mesin keadaan (\istilah{Finite State Machine}),
        dan pencarian jalur A*. Kecerdasan buatan berbasis pembelajaran mesin
        berada di luar cakupan.
  \item \textbf{Aset yang dipakai adalah aset siap pakai} yang boleh
        digunakan secara bebas, dengan penekanan pada cara mengimpor dan
        mengonfigurasinya. Pembuatan model dan animasi dari nol tidak dibahas.
  \item \textbf{Aspek bisnis dan distribusi permainan tidak dibahas}:
        monetisasi, penerbitan ke toko aplikasi, \istilah{multiplayer}
        melalui jaringan, dan penyelenggaraan peladen berada di luar lingkup
        mata kuliah ini.
\end{itemize}

\subsection{Capaian Akhir dan Bukti Belajarnya}

Empat capaian pembelajaran mata kuliah (\CPMK) yang harus Anda kuasai
dirumuskan pada Tabel~\ref{tab:capaian-bukti}. Perhatikan kolom terakhir:
setiap capaian ditopang oleh artefak nyata yang dapat diperiksa, bukan oleh
nilai ujian semata. Cara berpikir seperti ini akan Anda temui berulang kali
sepanjang buku dan diuraikan lebih jauh pada
Bagian~\ref{sec:kurikulum-obe}.

\begin{table}[htbp]
\centering
\caption{Capaian akhir mata kuliah dan bukti belajarnya}
\label{tab:capaian-bukti}
{\setlength{\tabcolsep}{4pt}\footnotesize%
\begin{tabularx}{\textwidth}{@{}C{1.7cm} L{4.6cm} L{3.6cm} Y@{}}
\toprule
\headrow \textbf{CPMK} & \textbf{Kemampuan akhir} & \textbf{Penopang} & \textbf{Bukti yang diperiksa} \\
\midrule
\CPMK{1} &
Merancang \istilah{game loop} dan logika fisika sederhana. &
Bab III--VII (Materi 1--5) &
Prototipe permainan 2D berbasis Pygame; tugas praktikum 01--05; jawaban UTS. \\
\addlinespace
\CPMK{2} &
Menerapkan algoritma kecerdasan buatan dasar untuk NPC. &
Bab VIII--XI (Materi 6--9) &
Mini proyek AI berisi NPC dengan mesin keadaan dan A*, beserta laporan
\istilah{playtest}. \\
\addlinespace
\CPMK{3} &
Menggunakan \istilah{game engine} Unity untuk membangun level. &
Bab XII--XIII (Materi 10--11) &
Level Unity yang dapat dimainkan beserta kontrol pemain dan konfigurasi
fisikanya. \\
\addlinespace
\CPMK{4} &
Mengintegrasikan aset audio dan visual ke dalam interaksi \istilah{game}. &
Bab XIV--XVI (Materi 12--14) &
Prototipe Unity lengkap dengan aset visual, animasi, antarmuka, dan audio;
laporan \istilah{playtest}; presentasi demo. \\
\bottomrule
\end{tabularx}}
\tabsource{Sumber: disusun dari RPS mata kuliah \NamaMK\ (\KodeMK) dan pemetaan
CPL--CPMK--Sub-CPMK.}
\end{table}
